% =====================================================================
\section{Đặt vấn đề và mục tiêu}
\label{sec:intro}
% =====================================================================

\subsection{Đặt vấn đề}
\label{sec:intro-problem}

Bài toán được khảo sát là phân loại ảnh CIFAR-10 vào hai siêu lớp do đề tài quy ước:
\emph{Animal} và \emph{Vehicle}. Mô hình dùng chung một bộ trích xuất đặc trưng ResNet-18 và
đồng thời học hai nhiệm vụ: nhiệm vụ chính là phân loại siêu lớp có giám sát, nhiệm vụ phụ là
dự đoán phép xoay đã được áp dụng lên ảnh. Nhãn của nhiệm vụ phụ được sinh tự động từ chính
phép biến đổi, nhờ đó không cần nhãn ngữ nghĩa của ảnh~\cite{he2016resnet,gidaris2018rotnet}.

Động cơ nghiên cứu không nằm ở việc đề xuất một kiến trúc mới, mà ở việc \emph{đánh giá khả
năng khai thác tín hiệu tự giám sát khi số ảnh có nhãn dùng cho huấn luyện bị giới hạn}.
Nhiệm vụ dự đoán góc xoay được vận dụng từ hướng nghiên cứu RotNet; việc ghép ResNet-18 với
hai đầu ra không được xem là đóng góp kiến trúc. Cần nhấn mạnh một điểm dễ gây nhầm lẫn:
\emph{góc xoay ở đây là góc biến đổi nhân tạo}, không phải góc tư thế thực của đối tượng
trong ảnh~\cite{gidaris2018rotnet}.

Ở cấu hình có sử dụng thêm ảnh không nhãn, thiết lập thực chất là \emph{học bán giám sát} với
một nhiệm vụ phụ tự giám sát. Thuật ngữ \enquote{học đa nhiệm} mô tả giai đoạn huấn luyện;
khi suy luận, hệ thống chỉ giữ nhiệm vụ phân loại Animal/Vehicle. Vì vậy lợi ích cần kiểm
chứng là \emph{chất lượng phân loại}, chứ không phải mặc nhiên là giảm chi phí suy luận so với
một mô hình đơn nhiệm có cùng backbone.

Ba câu hỏi đánh giá được dùng xuyên suốt báo cáo:

\begin{table}[htbp]
  \centering
  \small
  \begin{tabular}{@{}p{4.6cm} p{10.4cm}@{}}
    \toprule
    Câu hỏi & Nội dung đánh giá \\
    \midrule
    C1 --- Hiệu quả phân loại &
    M8 có thay đổi Macro-F1 và Accuracy so với baseline chỉ học phân loại siêu lớp trong cùng
    ngân sách nhãn hay không? \\
    \addlinespace
    C2 --- Vai trò của dữ liệu &
    Kết quả khác nhau thế nào khi nhánh rotation chỉ dùng ảnh thuộc tập có nhãn và khi được
    dùng toàn bộ tập train? \\
    \addlinespace
    C3 --- Chi phí thực hiện &
    Chi phí huấn luyện tăng bao nhiêu; mô hình một đầu ra khi suy luận có số tham số, dung
    lượng và độ trễ như thế nào trên máy đo xác định? \\
    \bottomrule
  \end{tabular}
  \caption{Ba câu hỏi đánh giá và nội dung tương ứng. Câu C1 và C2 được trả lời bằng khối
  thực nghiệm chính (\secref{sec:results-main}); câu C3 bằng các phép đo chi phí ở
  \secref{sec:results-cost}.}
  \label{tab:questions}
\end{table}

\subsection{Mục tiêu và câu hỏi nghiên cứu}
\label{sec:intro-goal}

\textbf{Mục tiêu tổng quát.} Xây dựng và đánh giá một quy trình thực nghiệm có thể tái lập để
xác định ảnh hưởng của nhiệm vụ dự đoán góc xoay đối với phân loại siêu lớp tại các mức
$10\,\%$, $20\,\%$ và $50\,\%$ nhãn huấn luyện.

\textbf{Mục tiêu mở rộng.} Đề cương giữ hướng Contrastive Learning/SupCon ở phần \emph{mở rộng
tùy chọn} (các cấu hình \cfg{E3}--\cfg{E6}). Báo cáo này triển khai đầy đủ hướng đó và báo cáo
kết quả riêng tại \secref{sec:results-ext}, sau khi phần bắt buộc đã hoàn tất.

\textbf{Giả thuyết cần kiểm chứng.} Nhiệm vụ phụ \emph{có thể} cải thiện kết quả, đặc biệt khi
ít nhãn. Đây là một giả thuyết khoa học, không phải chỉ tiêu cam kết. Báo cáo phải trình bày
đầy đủ cả trường hợp không cải thiện hoặc suy giảm; không có yêu cầu nào buộc kết quả của M8
phải cao hơn baseline.

\begin{warnbox}
\textbf{Nguyên tắc báo cáo trung thực.} Mọi kết quả trong báo cáo này — kể cả kết quả âm —
đều được giữ nguyên như khi chạy thực nghiệm. Không có số liệu minh họa nào được đặt vào vị
trí kết quả thực nghiệm, và không có lượt chạy nào bị loại bỏ sau khi đã xem kết quả test.
Mục tiêu không được điều chỉnh sau khi biết kết quả.
\end{warnbox}

\subsection{Phạm vi}
\label{sec:intro-scope}

Giữ nguyên CIFAR-10, ResNet-18 điều chỉnh cho ảnh $32 \times 32$ và thời gian dự kiến bốn
tuần. Phần bắt buộc tập trung vào baseline và M8; Contrastive Learning/SupCon là mở rộng tùy
chọn, được triển khai và báo cáo ở \secref{sec:results-ext} nhưng \emph{không} thuộc tiêu chí
hoàn thành bắt buộc. Việc đánh giá chi phí suy luận phục vụ thảo luận về Edge AI, \emph{không
thay thế} thử nghiệm trên thiết bị IoT thực tế và \emph{không mặc nhiên} chứng minh khả năng
triển khai trên vi điều khiển.

\begin{limbox}
\textbf{Những gì báo cáo này không khẳng định:} (i)~không so sánh trực tiếp Accuracy hai lớp
với bảng xếp hạng CIFAR-10 mười lớp; (ii)~không khẳng định đặc trưng học được tổng quát hơn
hay thuật toán là mới; (iii)~không khẳng định tối ưu năng lượng; (iv)~không khẳng định đã xây
dựng một hệ thống IoT hoàn chỉnh.
\end{limbox}
